% The epsilon-bias radius. Session 8 draft.
% Theory is settled (docs/R_EPSILON_THEORY.md, THEOREMS.md sections 16-24).
% Every number traces to a file under results/epsilon/. Do not edit one without the file.
% Convention reminder: r is the distance to the NEAREST failure, so the safe-move count is r-1.

\section{From a validity radius to a bias radius}
\label{sec:epsilon}

The breakdown radius says where the committed adjustment set stops being valid. It does not say how
wrong the reported number becomes once it does, and the two questions come apart: the shell profile
of the running example (Appendix~\ref{app:worked}) shows the bias rising gradually past $r_{\mathrm{val}}$
rather than jumping. An analyst told that their conclusion survives two revisions and fails at the
third still does not know whether the third costs them one percent or one hundred. This section
supplies the missing half, and it does so without weakening anything: the quantity defined here is
exactly zero on the ball $r_{\mathrm{val}}$ already certifies, so the validity radius is recovered as
the first point of a curve rather than replaced by it.

\subsection{The worst-case bias a knowledge state leaves open}

Fix the query $(X,Y)$, the set $Z$ read off $G_0$, and the observational law the analyst has, with
population covariance $\Sigma$. Write $\beta(S;\Sigma)$ for the coefficient on $X$ in the population
least-squares regression of $Y$ on $\{X\}\cup S$. The analyst reports a single number,
$\theta_Z=\beta(Z;\Sigma)$, which does not depend on which directed acyclic graph is the truth. What
the knowledge state leaves open is which one is, and each candidate $D$ assigns the effect its own
value $\tau_D=\beta(\mathrm{pa}_D(X);\Sigma)$, the quantity enumerated by \textsc{ida}
\citep{maathuis2009}, taken to be zero when $D$ makes $Y$ a parent of $X$. The ambiguity a state
leaves is therefore a finite set of real numbers, $T(G)=\{\tau_D : D\in[G]\}$, and the natural
diagnostic is how far the reported number can sit from any of them,
\begin{equation}
  B(G) \;=\; \max\{\,|\theta_Z-\tau| \;:\; \tau\in T(G)\,\}
        \;=\; \max\bigl(|\theta_Z-\min T(G)|,\;|\theta_Z-\max T(G)|\bigr).
  \label{eq:B}
\end{equation}
Two properties make \eqref{eq:B} the right object rather than one of the several alternatives. It is
an exact supremum over the ambiguity, not an average over resampled parameters, so it is an upper
bound on an error the analyst actually carries rather than on a hypothetical one; and because $T(G)$
is indexed by $[G]$, which grows as knowledge is withdrawn, it is monotone under the order on
$\mathfrak{G}_{\widehat{C}}$ by construction. A relative reading divides by $|\theta_Z|$, which is a
constant of the instance; dividing instead by anything that varies with the state --- the state's own
mean effect, say --- would read as a percentage and would destroy the monotonicity the certificate is
built on.

\begin{theorem}[Zero bias inside the certified shell]
\label{thm:zero}
If $d(G_0,G)\le r_{\mathrm{val}}-1$ then $T(G)=\{\theta_Z\}$, and hence $B(G)=0$.
\end{theorem}

The proof is one line: no state inside the shell fails, so $Z$ is a valid adjustment set in every
graph the state represents, and a valid adjustment set identifies the total effect. The content is in
how strong the conclusion is. It is not that the bias is small or insignificant; it is algebraically
zero, at the population level, for every observational law compatible with the equivalence class.
Inside the shell the only error left is sampling error, and there is nothing to recompute from state
to state because every state gives the same estimand. The converse fails in the benign direction:
at distance $r_{\mathrm{val}}$ some state admits a graph in which $Z$ is invalid, yet $\beta(Z;\Sigma)$
can still coincide with $\tau_D$ for the particular $\Sigma$ at hand, so $r_{\mathrm{val}}$ is a lower
bound on every bias radius defined below and understates robustness rather than overstating it.

\begin{theorem}[Monotonicity]
\label{thm:mono}
If $[G]\subseteq[H]$ then $T(G)\subseteq T(H)$ and $B(G)\le B(H)$. Consequently
$F_\varepsilon=\{G : B(G)>\varepsilon\}$ is upward-closed, and $F_{\varepsilon'}\subseteq F_\varepsilon$
whenever $\varepsilon\le\varepsilon'$.
\end{theorem}

Theorem~\ref{thm:mono} is monotonicity along the order, which is not the same as monotonicity along
the distance, and the difference is worth stating because the second reading is the tempting one and
it is false: two states at the same distance are typically incomparable, and a state one shell further
out can carry strictly less bias than a state nearer in. Section~\ref{sec:epsilon-eval} exhibits
minimal counterexamples. What is monotone in the distance is the worst case over the ball, which is
also the only thing a certificate could be built from, since a guarantee has to hold for whichever
wrong claims the analyst is in fact holding and not for an average one.

\subsection{The radius, and why it costs no more than the one we already compute}

Writing $\beta^{\uparrow}(d)$ for the largest $B$ over states within $d$ retractions of $G_0$, define
\begin{equation}
  r_\varepsilon \;=\; \min\{\,d(G_0,G) \;:\; B(G)>\varepsilon\,\} \;=\; \min\{\,d : \beta^{\uparrow}(d)>\varepsilon\,\}.
  \label{eq:reps}
\end{equation}
The reduction that makes $r_{\mathrm{val}}$ computable applies verbatim. Reading the proof of
Proposition~\ref{prop:monotone-reach} back, the only property of failure it consumes is that failure
is upward-closed; nothing about adjustment sets or $d$-separation enters. It is therefore a statement
about arbitrary up-sets, and Theorem~\ref{thm:mono} puts $F_\varepsilon$ among them, so the nearest
$\varepsilon$-violating state is again reachable by retraction alone and the search is again confined
to withdrawals of the analyst's own commitments. The up-set of $G_0$ is parameterised by the subsets
of $\mathcal{K}_{G_0}$, and shell $d$ is generated directly from the $d$-subsets, so a shell can be
built without traversing any shell beneath it.

Four consequences give the algorithm. Shells below $r_{\mathrm{val}}$ are never evaluated, by
Theorem~\ref{thm:zero}, so the expensive part of the computation is skipped on exactly the region the
certificate is strongest; the maximum over a ball is attained on its outermost sphere, so interior
states need no evaluation either; one traversal answers every $\varepsilon$, including one chosen
after the fact, because each radius is then a lookup on a monotone staircase rather than a separate
search; and monotonicity licenses bisection, locating the first crossing in a logarithmic number of
shells when the crossing is far out. Evaluating $B$ itself needs no graph enumeration: the ambiguity
set is determined by the possible parent sets of the treatment, and those are obtained by orienting
the undirected edges at $X$ each way and keeping the assignments whose Meek closure does not fail,
which is $2^{\deg(X)}$ polynomial closures where $\deg(X)$ counts undirected edges at the treatment
alone. Two-sided bounds are available without any exhaustive pass: $r_{\mathrm{val}}$ and any
smaller-$\varepsilon$ radius bound $r_\varepsilon$ from below at no cost, a single greedy retraction
chain that reaches a violating state bounds it from above with a replayable witness, and one
evaluation at $\widehat{C}$ --- the maximum of the order, hence the maximum of $B$ --- settles
$r_\varepsilon=\infty$ for every threshold above it.

\begin{proposition}[The certificate]
\label{prop:cert}
Suppose at most $k$ of the orientations of $\mathcal{K}_{G_0}$ are false of the truth. Then the
realised identification error obeys $|\theta_Z-\tau^{*}|\le\beta^{\uparrow}(k)$. In particular
$\theta_Z=\tau^{*}$ exactly when $k<r_{\mathrm{val}}$, the error is at most $\varepsilon$ when
$k<r_\varepsilon$, and when $r_{\varepsilon_x}\le k<r_{\varepsilon_y}$ the worst case lies in
$(\varepsilon_x,\varepsilon_y]$.
\end{proposition}

\begin{proposition}[The band is attained]
\label{thm:tight}
For every $k$ there is a graph $D\in[\widehat{C}]$ that fits the observed law exactly, contradicts at
most $k$ of the analyst's orientations, and realises $|\theta_Z-\tau_D|=\beta^{\uparrow}(k)$.
\end{proposition}

Take a state at distance at most $k$ attaining $\beta^{\uparrow}(k)$ and, inside it, a graph attaining
the maximum. It is Markov equivalent to the truth, so it fits $\Sigma$; and it withholds at most $k$
of $\mathcal{K}_{G_0}$, by the parameterisation of the up-set. So the bound is realised by a world the
analyst cannot rule out on the evidence they hold, and the band cannot be tightened without an
assumption beyond the CPDAG, the data and the error budget. Refining the $\varepsilon$ grid refines
the band; nothing else does.

The proof is that the state obtained by withdrawing exactly the false orientations still represents
the true graph, so the truth is inside the ambiguity set that state leaves and the bound applies to
it. The budget is denominated in orientations of the closure, the same unit as $r_{\mathrm{val}}$, and
not in asserted claims: Meek's rules cascade, so a small number of wrong assertions can produce a
larger number of wrong closure orientations, and reading the budget in the wrong unit makes a sound
certificate look violated. The claim-level statement is the analogue built on $r_{\mathrm{claim}}$.

\subsection{What this changes about the assumption the radius carried}
\label{sec:epsilon-assumption}

Everything above rests on the same order-theoretic chain as $r_{\mathrm{val}}$ and adds nothing to it,
so it is worth recording that the chain is now complete. The anti-exchange property, on which the
retraction reduction depends, had been verified rather than proved in the case of two distinct
undirected edges, and every radius was accordingly reported with a one-sided caveat. The obstruction
was that the classical argument reorients freely inside a chordal chain component, and an MPDAG
carrying background knowledge need not have one. The missing step is supplied instead by the form of
Chickering's transformational theorem that counts its reversals: between Markov equivalent graphs
differing on $m$ edges there is a sequence of exactly $m$ distinct covered-edge reversals, and a walk
of $m$ unit steps taking the difference count from $m$ to zero must decrease at every step, so every
reversal flips an edge on which the two graphs currently differ and no edge on which they agree is
ever touched. Two extensions of the same MPDAG agree on all of its directed edges, so every
intermediate graph retains them and stays inside the class the MPDAG represents, which is exactly the
premise the anti-exchange argument needed. Appendix~\ref{app:order} carries the corrected proof.
Anti-exchange, and with it the covering lemma, upper semimodularity, the non-expansion of distance
under the join and monotone reachability, are therefore unconditional, and the caveat attached to
every reported radius in the preceding sections can be dropped. No previously reported number
changes; what changes is that ``exact if the conjecture holds'' becomes ``exact''.

\subsection{Evaluation}
\label{sec:epsilon-eval}

Everything below is traced to a file under \texttt{results/epsilon/}; nothing here is estimated from
memory. Tolerances are absolute and stated; no tolerance was widened to make a check pass.

\paragraph{The predictions the construction has to survive.}
Seven consequences of Section~\ref{sec:epsilon} are falsifiable by direct computation, and were
checked over 400 screened instances crossed with three independent parameter draws each
(\texttt{results/epsilon/verification/}): that the bias vanishes identically inside the certified
shell (1{,}758 states); that it is monotone over every ordered pair of the enumerated space
(372{,}099 pairs); that the staircase is non-decreasing and attained on its sphere (17{,}160 shells);
that the incremental, bisection and brute-force radii coincide (5{,}844 comparisons); that the
radius is non-decreasing in the threshold and never below $r_{\mathrm{val}}$ (5{,}844); that the
realised error respects the certificate at the budget the closure actually denominates (22{,}680);
and that the semi-local ambiguity set equals the enumerated one (11{,}413). Across
\textbf{436{,}798 checks there were no violations}, at a tolerance of $10^{-9}$.

\paragraph{The staircase, on the running example.}
Figure~\ref{fig:staircase} replaces the sampled-mean column of Appendix~\ref{app:worked} with the
exact worst case. The two disagree in shape, and the disagreement is informative. The mean rises
shell by shell; the worst case, at the documented parameter draw, is flat at zero out to
$r_{\mathrm{val}}$ and then jumps straight to $B(\widehat{C})$ --- the largest bias any perturbation
anywhere can produce --- and stays there. What rises gradually with depth is the *proportion* of
states at which $Z$ fails, not how badly the worst one fails, so the apparent ramp is an artefact of
averaging over an increasingly contaminated shell. Over 400 further parameter draws the shape splits
roughly evenly (221 flat, 179 with one further rise), so neither shape is universal; what is stable
is that on this instance the entire bias risk arrives at $r_{\mathrm{val}}$. The radii computed over
the retraction up-set agree with a brute-force breadth-first search over the whole enumerated space
in all three knowledge states, which is Theorem~\ref{thm:mono} and its corollaries holding on a case
where the two code paths share nothing but the definition.

Two further readings matter. Depths beyond $|\mathcal{K}_{G_0}|$ carry no entry, and that is not a
gap: an analyst holding $|\mathcal{K}_{G_0}|$ orientations cannot have more than that many of them
wrong, so those depths lie outside Proposition~\ref{prop:cert} by construction. And
$r_\varepsilon = r_{\mathrm{val}}$ is not a null result: even where tolerating $\varepsilon$ buys no
extra certified shell, it supplies the magnitude that $r_{\mathrm{val}}$ cannot. The validity radius
says the third revision breaks the set; the band says what breaking it costs.

\paragraph{Two naive readings, refuted.}
The construction maximises over a shell and uses a supremum rather than a mean, and both choices are
forced (\texttt{results/epsilon/counterexamples/}).

\emph{Bias is not monotone in distance.} At four nodes --- the smallest size at which the question is
non-trivial, and where 355 witnesses exist --- there is an instance with a state at distance one
carrying bias $1.34$ and a state at distance two carrying exactly zero. The two are incomparable in
the order, so Theorem~\ref{thm:mono} is untouched; what fails is the reading of it that confuses
distance with knowledge. A certificate read off a single perturbation path would be wrong here,
which is why $\beta^{\uparrow}$ maximises over the shell.

\emph{The mean-based alternative is not merely unjustified; it is wrong.} On identical pairs, a
sampled mean absolute bias is non-monotone under the order on at least $34.7\%$ of comparable pairs
($278$ of $801$, counting a pair only when the deficit exceeds three combined Monte-Carlo standard
errors in each of two independent repetitions, so the rate is a lower bound; median margin $9.6$
standard errors), while the worst case defined here is non-monotone on \textbf{none} of them. The
consequence is operational: because $\{$mean $> \varepsilon\}$ is not an up-set, the retraction-only
search is not licensed for it, and comparing it against brute force over the whole space
disagrees on $269$ of $644$ threshold probes --- in the observed cases by reporting that no
perturbation reaches the threshold anywhere when one does at distance two. That is not a
conservative approximation but an answer in the direction that overstates robustness.

\paragraph{What the refinement buys, over 900 instances.}
The sweep (\texttt{results/epsilon/study/}, 900 accepted instances over five generators, node counts
6--12, three knowledge coverages and four corruption regimes) says the payoff is not the one the
framing suggests. Extra certified shells are rare: $r_\varepsilon$ exceeds $r_{\mathrm{val}}$ on
$1.0\%$ of instances at $\varepsilon = 1\%$, rising to $24.5\%$ at $25\%$ and $75.6\%$ at $100\%$,
and when it does the gain is almost always exactly one shell (mean $1.00$--$1.14$). What rises much
faster is the fraction of instances where the threshold is \emph{never} reached: from $1.4\%$ at
$\varepsilon = 1\%$ to $70.1\%$ at $\varepsilon = 100\%$. The practical content of the
$\varepsilon$-radius is therefore a \emph{ceiling} rather than a longer runway --- on seven instances
in ten, no perturbation of the analyst's knowledge, however large, can move the estimate by its own
magnitude. That is a statement $r_{\mathrm{val}}$ cannot make at all, and it is available from the
same traversal.

Three further readings. First, $r_0 = r_{\mathrm{val}}$ on all 900 instances with no exceptions,
which is the genericity in the converse to Theorem~\ref{thm:zero} behaving as expected: coincidental
cancellation at the breakdown radius is possible but was never observed. Second, the staircase is a
step in $84.1\%$ of instances and a ramp in $15.1\%$, so the worst case usually saturates the moment
validity fails --- the gradual rise visible in a sampled mean is, in the large majority of cases, an
artefact of averaging over a shell whose contaminated fraction is growing. Third, the certificate
held on every one of the 884 instances where $\beta^{\uparrow}(k)$ was actually computed, with
\textbf{zero violations}; the remaining 16 are instances where the traversal stopped before depth
$k$, so the recorded value is a lower bound and the comparison is not a test of
Proposition~\ref{prop:cert} (recomputing their full staircases, the certificate holds there too).
Among instances with a non-degenerate bound the conservativeness ratio has \emph{median} $1.000$ and
maximum $1.000$: the bound is not merely valid but routinely attained, which is
Theorem~\ref{thm:tight} in the data.

\paragraph{Cost.}
Both exact strategies agree on every instance (zero disagreements). The incremental traversal
evaluates a median of one state and a mean of $2.09$; bisection a median of one and a mean of
$5.14$, and up to $172$ against the incremental maximum of $65$. The reason is that
$r_\varepsilon$ is small almost everywhere, so a probe near the middle of the subset lattice is
wasted work; bisection is the right strategy only in the regime this corpus does not contain, and
the dispatch --- incremental under a shell budget, bisection as a fallback --- fires its fallback
rarely. The anytime greedy chain, which enumerates no shell at all, returns the exact radius on
$98.4\%$ of instances, so an upper bound is usually free. Skipping the certified shells saves a
median of one and up to six bias evaluations per instance on the measured subsample; the saving is
modest here because $r_{\mathrm{val}}$ is itself usually one, and would grow with it.

\paragraph{Against the incumbent.}
On a matched subsample, the radius defined here disagrees with the mean-based quantity currently in
the codebase on at least one threshold in $84\%$ of instances, in both directions (23 larger, 14
smaller, 8 equal, 5 mixed). The disagreement is not a tie-break between two reasonable choices: on
the same instances the mean was non-monotone under the order in 49 of 50 cases, so the search used
to compute its radius is not licensed there at all.

\paragraph{Finite samples.}
$r_{\mathrm{val}}$ reads no numbers and has no sampling distribution; $r_\varepsilon$ reads a
covariance and does. Over 120 instances at sample sizes $100$ to $10{,}000$
(\texttt{results/epsilon/finite\_sample/}), the plug-in radius is too large --- the direction that
overstates robustness --- in up to $5.0\%$ of cases at $n=100$, falling to $0$--$0.8\%$ at
$n=10{,}000$. Thresholding a bootstrap upper bound on the bias instead, and reporting the $5\%$
quantile of the resulting radius, drives that rate to \textbf{$0.0\%$ in every cell}, at the cost of
understating the radius on $1$--$17\%$ of instances. The conservative radius is the one to report.
